\ifx\allfiles\undefined
\documentclass[12pt, a4paper, oneside, UTF8]{ctexbook}
\def\path{../config}
\input{\path/package.tex}

% 定理环境设置
\input{\path/theorem1_zh.tex}

\input{\path/custom.tex}

% 修改参数改变封面样式，0 默认原始封面、内置其他1、2、3种封面样式
\def\myIndex{3}


\ifnum\myIndex>0
    \input{\path/cover_package_\myIndex} 
\fi

\def\myTitle{复变函数与积分变换笔记}
\def\myAuthor{M.K.}
\def\myDateCover{\today}
\def\myDateForeword{\today}
\def\myForeword{前言}
\def\myForewordText{
    
    这是一个基于\LaTeX{}模板的复变函数与积分变换笔记．
    内容参考了包括但不限于以下资料：
    \begin{itemize}
        \item 《复变函数与积分变换》（科学出版社）第三版
        \item 《复变函数习题精选精讲》（山东科学技术出版社）第一版
    \end{itemize}

    封面来源于\href{https://www.pixiv.net/users/2642047}{アシマ / Ashima}的作品\href{https://www.pixiv.net/artworks/147092236}{Assembly}
    如有侵权请联系我删除．
}
\def\mySubheading{副标题}


\begin{document}
%\input{../config/cover}
\else
\fi
\chapter{解析函数}
\section{复变函数微分}

\subsection{导数}
\begin{definition}[导数]
设 $D\subseteq\C$ 为非空开集，$f\colon D\to\C$，$z_0\in D$。若极限
\[
f'(z_0)\defeq\lim_{z\to z_0}\frac{f(z)-f(z_0)}{z-z_0}
\]
存在，则称 $f$ 在 $z_0$ 处\textbf{可导}，称 $f'(z_0)$ 为 $f$ 在 $z_0$ 处的\textbf{导数}。
\end{definition}

\begin{definition}[导函数]
若 $f$ 在 $D$ 内每一点均可导，则称 $f$ 在 $D$ 内\textbf{可导}，并称
\[
f'(z) \defeq \lim_{\Delta z\to 0}\frac{f(z+\Delta z)-f(z)}{\Delta z},\quad z\in D
\]
为 $f$ 的\textbf{导函数}。
\end{definition}

\subsection{解析}
\begin{definition}[解析]
设 $D\subseteq\C$ 为非空开集，$f\colon D\to\C$。若 $f$ 在 $z_0\in D$ 的某邻域内处处可导，则称 $f$ 在 $z_0$ 处\textbf{解析}（analytic / holomorphic）；若 $f$ 在 $D$ 内每一点均解析，则称 $f$ 在 $D$ 内\textbf{解析}，或称 $f$ 是 $D$ 内的\textbf{解析函数}。
\end{definition}

\begin{theorem}[Cauchy--Riemann 方程（必要条件）]
设 $D\subseteq\C$ 为开集，$f(z)=u(x,y)+iv(x,y)$ 在 $z_0=x_0+iy_0\in D$ 处可导，且 $u,v$ 在 $(x_0,y_0)$ 处可微，则
\[
\frac{\partial u}{\partial x}(x_0,y_0) = \frac{\partial v}{\partial y}(x_0,y_0),\qquad
\frac{\partial u}{\partial y}(x_0,y_0) = -\frac{\partial v}{\partial x}(x_0,y_0).
\]
\end{theorem}

\begin{theorem}[Cauchy--Riemann 方程（充要条件）]
设 $D\subseteq\C$ 为开集，$f(z)=u(x,y)+iv(x,y)$。若 $u,v$ 在 $(x_0,y_0)$ 的某邻域内可微，且满足 Cauchy--Riemann 方程
\[
\frac{\partial u}{\partial x} = \frac{\partial v}{\partial y},\quad \frac{\partial u}{\partial y} = -\frac{\partial v}{\partial x},
\]
则 $f$ 在 $z_0$ 处可导，且
\[
f'(z_0) = \frac{\partial u}{\partial x}(x_0,y_0) + i\frac{\partial v}{\partial x}(x_0,y_0)
= \frac{\partial v}{\partial y}(x_0,y_0) - i\frac{\partial u}{\partial y}(x_0,y_0).
\]
\end{theorem}

\begin{remark}
若 $u,v$ 在 $D$ 内连续可微，且 Cauchy--Riemann 方程在 $D$ 内处处成立，则 $f$ 在 $D$ 内解析。
\end{remark}

\section{初等函数}
本节约定 $z=x+iy$，$x,y\in\R$。

\subsection{指数函数 $\e^z$}
\begin{definition}[指数函数]
${\e}^z\colon\C\to\C$ 定义为
\[
\e^z \defeq \sum_{n=0}^{\infty}\frac{z^n}{n!}.
\]
\end{definition}

\begin{theorem}[指数函数的性质]
\begin{enumerate}
    \item $\e^{z_1+z_2} = \e^{z_1}\cdot \e^{z_2}$;
    \item $\e^z \neq 0$;
    \item $\overline{\e^z} = \e^{\bar{z}}$;
    \item $\e^{z+2k\pi i}=\e^z$，$k\in\Z$；
    \item $\abs{\e^z} = \e^x$。
\end{enumerate}
\end{theorem}

\subsection{对数函数 $\Log z$}
\begin{definition}[对数函数（多值）]
设 $z\in\C\setminus\{0\}$。$z$ 的\textbf{对数}（多值）定义为满足
\[
\e^w = z
\]
的所有 $w\in\C$。记作 $\Log z$。
\end{definition}

\begin{theorem}[对数函数（多值）公式]
对任意 $z = r\e^{i\theta}$（$r>0$，$\theta\in\R$），
\[
\Log z = \ln r + i(\theta + 2k\pi) = \ln r + i\Arg z,\quad k\in\Z.
\]
\end{theorem}

\begin{definition}[对数主值 $\log z$]
取 $\arg z\in(-\pi,\pi]$，定义
\[
\log z \defeq \ln\abs{z} + i\arg z.
\]
称为\textbf{对数主值}，它在 $\C\setminus(-\infty,0]$ 上单值解析，且 $(\log z)' = \dfrac{1}{z}$。
\end{definition}

\subsection{幂函数 $z^\alpha$}
\begin{definition}[一般幂函数]
设 $\alpha\in\C$，$z\in\C\setminus\{0\}$。定义
\[
z^\alpha \defeq \e^{\alpha\log z} = \e^{\alpha(\ln\abs{z} + i\arg z)}.
\]
\end{definition}

\begin{theorem}[一般幂函数的性质]
当 $\alpha\in\R$ 时，$z^\alpha$ 的模为 $\abs{z}^\alpha$，辐角为 $\alpha\arg z$；特别地 $\alpha=\frac{1}{n}$（$n\in\N$）时，$z^{\frac{1}{n}}$ 为 $z$ 的 $n$ 次方根。
\end{theorem}

\subsection{三角函数}
\begin{definition}[正弦与余弦]
对任意 $z\in\C$，
\[
\sin z \defeq \frac{\e^{iz} - \e^{-iz}}{2i},\qquad
\cos z \defeq \frac{\e^{iz} + \e^{-iz}}{2}.
\]
\end{definition}

\begin{theorem}[三角函数的性质]

$$\overline{\sin z} = \sin \bar{z},\quad \overline{\cos z}=\cos\bar z$$

\end{theorem}

\begin{definition}[正切与余切]
\[
\tan z \defeq \frac{\sin z}{\cos z},\quad \cot z \defeq \frac{\cos z}{\sin z}.
\]
\end{definition}

\subsection{双曲函数}
\begin{definition}[双曲正弦与余弦]
\[
\sinh z \defeq \frac{\e^z - \e^{-z}}{2},\qquad
\cosh z \defeq \frac{\e^z + \e^{-z}}{2}.
\]
\end{definition}

\begin{theorem}[双曲函数的性质]
\begin{enumerate}
    \item $\cosh^2 z - \sinh^2 z = 1$。
    \item $\sinh(-z)=-\sinh z$（奇），$\cosh(-z)=\cosh z$（偶）。
    \item $\overline{\sinh z}=\sinh\bar z$，$\overline{\cosh z}=\cosh\bar z$。
    \item $(\sinh z)' = \cosh z$，$(\cosh z)' = \sinh z$。
    \item 与三角函数的关系：$\sinh z = -i\sin(iz)$，$\cosh z = \cos(iz)$。
\end{enumerate}
\end{theorem}

\begin{definition}[双曲正切与余切]
\[
\tanh z \defeq \frac{\sinh z}{\cosh z},\quad \coth z \defeq \frac{\cosh z}{\sinh z}.
\]
\end{definition}

\section{调和函数}
\subsection{调和函数的定义}
\begin{definition}[调和函数]
设 $D\subseteq\R^2$ 为开集，$u\colon D\to\R$。若 $u$ 在 $D$ 内二阶连续可微，且满足 \textbf{Laplace 方程}
\[
\Delta u \defeq \frac{\partial^2 u}{\partial x^2} + \frac{\partial^2 u}{\partial y^2} = 0,
\]
则称 $u$ 为 $D$ 内的\textbf{调和函数}。
\end{definition}

\begin{theorem}[解析函数的实部与虚部皆调和]
设 $D\subseteq\C$ 为开集，$f=u+iv$ 在 $D$ 内解析。则 $u,v$ 均为 $D$ 内的调和函数。
\end{theorem}

\subsection{共轭调和函数}
\begin{definition}[共轭调和函数]
设 $D\subseteq\C$ 为开集，$u\colon D\to\R$ 调和。若存在 $v\colon D\to\R$ 使得 $f=u+iv$ 在 $D$ 内解析，则称 $v$ 为 $u$ 在 $D$ 内的\textbf{共轭调和函数}。
\end{definition}

\begin{theorem}[Cauchy--Riemann 方程的调和等价形式]
设 $u,v\in C^2(D,\R)$。$f=u+iv$ 在 $D$ 内解析当且仅当 $u,v$ 均为 $D$ 内的调和函数，且满足 Cauchy--Riemann 方程
\[
\frac{\partial u}{\partial x} = \frac{\partial v}{\partial y},\qquad \frac{\partial u}{\partial y} = -\frac{\partial v}{\partial x}.
\]
\end{theorem}

\ifx\allfiles\undefined
\end{document}
\fi